\section{Implementation}
\label{sec:impl}

The system is written in Python (about 79 commits, 818 unit tests, all passing at the time of writing). It is organised as a \emph{harness of four plugin packages} around a SQLite store, with the three knowledge bases beside it (Figure~\ref{fig:harness}). Each package root holds only a contract (the data a stage takes and returns) and a registry (versions by name); every method-specific piece lives in a version folder, so a stage can be replaced by registering a new version without touching its neighbours. This mattered in practice: the change of direction of Section~\ref{sec:idea-pivot} was carried out by adding new versions of three stages while the old ones stayed selectable as controls.

\begin{figure}[t]
\centering
\fitwidth{\begin{tikzpicture}[
  >=Stealth, font=\scriptsize,
  stage/.style={draw,thick,rounded corners=3pt,fill=cLight,align=center,minimum height=13mm,text width=27mm,inner sep=3pt},
  pkg/.style={font=\scriptsize\bfseries},
  ver/.style={font=\tiny,gray},
  store/.style={draw=cStore,thick,fill=cStore!12,rounded corners=3pt,align=center,text width=30mm,inner sep=3pt,minimum height=11mm},
  kb/.style={draw,thick,rounded corners=3pt,align=center,text width=26mm,inner sep=3pt,minimum height=11mm},
  arr/.style={->,thick,gray!80},
  io/.style={draw=cDoc,thick,fill=cDoc!12,rounded corners=6pt,align=center,inner sep=4pt}
]
% ---------- ingest path (top row) ----------
\node[io] (doc) at (0,0) {a filing\\(PDF, HTML, markdown)};
\node[stage,right=8mm of doc] (da) {\textbf{1. document adapters}\\ \code{parse(path) $\to$ ParsedDocument}\\ \textcolor{gray}{mineru\_pdf $\cdot$ plain\_text/html $\cdot$ docfinqa\_markdown}};
\node[stage,right=8mm of da] (ing) {\textbf{2. ingestion}\\ \code{ingest(doc) $\to$ report}\\ \textcolor{gray}{v1 facts into the graph $\cdot$ \textbf{v2} text + tables, no facts}};
\node[store,right=8mm of ing] (st) {\textbf{document store} (SQLite)\\ chunks + vectors, sentences,\\ headings $\to$ parts, documents (cover profile)};
\draw[arr] (doc) -- (da); \draw[arr] (da) -- (ing); \draw[arr] (ing) -- (st);

% ---------- query path (bottom row) ----------
\node[io,below=22mm of doc] (q) {a question\\(+ optional document,\\ answer format)};
\node[stage,right=8mm of q] (loop) {\textbf{4. agent loops}\\ \code{answer(question) $\to$ AgentAnswer}\\ \textcolor{gray}{v2 graph-as-store $\cdot$ v3 no graph $\cdot$ \textbf{v4} graph as prior}};
\node[stage,right=8mm of loop] (ret) {\textbf{3. retrieval}\\ \code{retrieve(need) $\to$ passages}\\ \textcolor{gray}{v1 solver + text $\cdot$ v2 ladder $\cdot$ \textbf{v3} learned scorer}};
\node[io,below=10mm of loop] (ans) {answer: value or text, cited cells,\\ derivation, assumptions, confidence};
\draw[arr] (q) -- (loop); \draw[arr] (loop) -- (ret); \draw[arr] (ret.north) -- ++(0,5mm) -| (st.south);
\draw[arr] (loop) -- (ans);

% ---------- knowledge bases (right column) ----------
\node[kb,draw=cEntity,fill=cEntity!10,right=10mm of st,yshift=0mm] (cg) {\textbf{concept graph}\\ entities, equations,\\ subtotals, rules, links};
\node[kb,draw=cDoc,fill=cDoc!10,below=3mm of cg] (sg) {\textbf{source graph}\\ kinds, parts, periods};
\node[kb,draw=cConnector,fill=cConnector!12,below=3mm of sg] (br) {\textbf{bridge}\\ where printed, how often};
\node[kb,draw=cStore,fill=cStore!10,below=3mm of br] (co) {\textbf{company store}\\ filers, names over time,\\ co-registrants};
\begin{scope}[on background layer]
\node[draw=gray,dashed,rounded corners=4pt,fit=(cg)(sg)(br)(co),inner sep=3mm,label={[gray,font=\scriptsize]above:knowledge bases (\code{knowledge/})}] (kbs) {};
\end{scope}
\draw[arr,dashed] (kbs.west |- loop) -- (ret.east |- loop) ;

% ---------- caption inside: contracts ----------
\node[gray,font=\scriptsize,align=left,below=4mm of ans.south west,anchor=north west] {Each numbered package root holds a \emph{contract} and a \emph{registry} only; versions are plugins selected by name.\\ Dashed: the knowledge bases are read by the adapters (candidate kinds for the cover), the retriever (the bridge's share, the parts) and the loop (formulas, chances, documents).};
\end{tikzpicture}}
\caption{The harness. Top: the ingest path turns a filing into stored text, tables, headings and a cover profile. Bottom: the query path answers a question through the agent loop, which draws passages from the retriever and formulas, chances and document knowledge from the knowledge bases. Bold version names are the registry defaults.}
\label{fig:harness}
\end{figure}


\subsection{The data model and store}
\label{sec:impl-model}

Entities, connectors and edges are plain Python dataclasses; the symbolic bodies are sympy expressions~\citep{meurer2017sympy}. A \code{NodeRef} tags a reference as entity or connector, which is what allows a binding to point at another connector. Connector kinds are registered by a decorator, so a new kind (rule, aggregation) plugs in with no schema change. The store is SQLite: entities, connectors with their serialised expression and bindings, edges, one table of float32 vectors keyed by node and embedding model, and the extension tables for documents (chunks, sentences, headings, document profiles) added during the year. Connectors are loaded in topological order of their connector-to-connector dependencies. A survey of the knowledge-graph management literature in June confirmed that at this scale (thousands of nodes, tens of thousands of edges) an in-memory adjacency structure over a SQLite file is the endorsed design, and that entity resolution, not storage, is the scaling risk; that proved exactly right.

\subsection{Building the concept graph}
\label{sec:impl-concept}

\paragraph{Textbook equations.}
The first source was finance textbooks. PDFs are parsed with MinerU~\citep{wang2024mineru}, a vision-language layout model that returns formulas as LaTeX and tables as structured JSON; a section-aware chunker keeps each equation with the text around it; one model call per equation returns the canonical \(\ldots=0\) expression, the bound entities with definitions, units and types, applicability conditions and a plain-English description. Entities are de-duplicated by exact name and then by embedding similarity; two equations are merged only if their bodies agree \emph{and} they bind the same entity set. A hand-curated corpus of 22 domains (about 2,600 equations, 1,300 rules and 150 facts, written in the pipeline's output shape) was ingested through the same builders. The result, \code{finance.db}, held 4,208 entities and 2,502 equations by July.

\paragraph{The SEC concept graph, verified against filings.}
\label{sec:impl-secgraph}
The second source is the US reporting taxonomy: 9,330 concepts (US GAAP and IFRS) with definitions and, from 802,000 filer captions, the labels filers actually print for each. These concepts had no relations. Rather than let a model invent them, \emph{every candidate relation is measured before admission}: the SEC's quarterly numeric data (4.6 million consolidated values over five quarters) is the answer key, a candidate is evaluated over every annual filing that reports all its operands at 0.5\% tolerance, and the builder refuses a rate below 0.5, demotes 0.5--0.7 to ``empirical'' and admits the rest as identities. Three families were built:
\begin{itemize}
\item \emph{Subtotals are hierarchy, not formula.} Each filer sums a different subset of lines, so a subtotal is an \emph{aggregation connector} with a signed, weighted part edge per member ever observed. Members come from the taxonomy's calculation linkbases and from each filer's own arithmetic, recovered from the presentation order by a stack reduction that accepts a block only because its signed values sum to the total (221,000 blocks; 91--97\% of annual filings explained for the main balance-sheet and cash-flow totals).
\item \emph{Roll-forwards} (the balance each ``increase/decrease in'' concept moves) chosen by measurement, admitted as empirical because working-capital roll-forwards hold only 35--60\% of the time.
\item \emph{Textbook laws rebound onto concepts}: 1,292 textbook equations offered to authoring agents with embedding shortlists; 130 relations admitted after independent review; one harmonisation pass over the family.
\end{itemize}
An IFRS mirror (206 aggregations, 24 identities) and a reviewed IFRS--US GAAP bridge of 292 same-quantity pairs followed. The verification exposed defects a reader would not have caught: cash-flow totals carry no debit/credit type in the taxonomy; shell companies with equity near zero create coincidental sums; a textbook ``pre-tax income'' identity held in only 36\% of filings and was refused.

\paragraph{Repairs, merges and the single graph.}
The textbook graph had three known defects, all found by measurement: statistical operators that made 14\% of equations unsolvable, entities defined wrongly or mixing two meanings, and one quantity split across several entities. Each was repaired by an \emph{authored decision file applied by a script that refuses the whole file on any decision that does not check out}, so nothing is edited by hand in the database:
\begin{itemize}
\item \emph{Statistics lifted} (September): 317 equations rewritten, 23 removed as mere definitions, 212 statistic entities created under a naming rule or an authored decision, the normal distribution function rewritten into sympy's \code{erf} so it evaluates and inverts.
\item \emph{Textbook--SEC links} (October): 1,188 candidate entities reviewed by twenty independent reviewers under one test (``a number filed under the concept is a value of the textbook quantity with nothing added, removed or re-based''), three adversarial reviewers, then a final read; 115 links kept, 54 rejected, 18 held back. The review also found 69 entities defined wrongly in the textbook graph (``total assets'' defined as a pension plan's assets; ``current assets'' as the balance-of-payments current account).
\item \emph{Repairs and merges} (October): those 69 entities redefined or split with every bound equation re-read; 18 duplicate entities merged into 11; 74 duplicate equations removed; the 2,058 SEC relations imported so that a concept linked to a textbook entity \emph{is} that entity.
\end{itemize}
The result is one graph, \code{finance\_sec.db}: 9,121 entities, 5,457 relations, 4,983 concept links with one entity per concept, a giant component of 6,860 nodes (from 2,948 in the textbook graph alone), and a vector for every entity. ``Total assets'' now carries 40 formulas and 5 filing-tested subtotals; ``net income'' 49 and 16.

\subsection{The solver}
\label{sec:impl-solver}

The solver is iterative-deepening AND--OR search with conditional memoisation. The outer loop is over \emph{trust layers} (identities only, then laws, then empirical relations); the inner loop deepens a cost budget, cost being the number of derived quantities in the tree, so the first derivation to close is the one of highest trust and least data demand. Failures are memoised with the budget they failed under and the goals they cycled on, so a shared dead end costs one exploration rather than one per route (the earlier uncached depth-first search ran for minutes on failure). Symbolic solves are cached and run without simplification, guarded by a numeric residual check. An aggregation gives its total as the weighted sum of the parts present at the coordinate and a part as the residual. Where a stored cell holds several distinct readings, a model is shown the readings with their provenance and chooses one or declines; a decline is not a failure but sends the goal through its connectors as for an empty cell.

\subsection{Reading documents}
\label{sec:impl-reading}

\paragraph{Parsing.}
PDFs go through MinerU; its throughput was raised about eightfold by running the layout model under vLLM across the lab's GPUs (the 20-page sample fell from 393~s to 44~s), and a batch driver parses a folder across cards resumably. Tables are converted to the pipe-separated grid every later stage reads, with the caption emitted on the line before the grid; code listings, equations and image captions are kept, a model's \emph{description} of an image is not, so nothing unprinted enters the store. EDGAR HTML filings go through a reader written for the retrieval corpus: a short block printed wholly in bold is a heading, tables become the same grids, every inline-XBRL fact is recorded with its context.

\paragraph{Cover and outline.}
Two model calls per document establish what it is. The \emph{cover reading} shows the first three pages and the candidate kinds the source graph offers, narrowed in tiers (an identifying phrase that titles a cover line; a weak phrase; the named filer's class and the dates printed; nearest in meaning), and returns the filer, every registrant on a combined filing, the kind, the period end, the parts printed on the cover and a description; every value is checked against the printed text (a quoted title must be a title line, not words inside a sentence) and a company name not printed is dropped. The \emph{outline} shows the model the headings and nothing else and asks for the same headings indented; the reply is checked line by line (nothing dropped, added, reworded or reordered), so position identifies each heading and indentation gives the nesting the layout model never reports. Headings are then linked to the source graph's parts (``Item 8'' and its children), and the kind's skeleton is shown as a hint for depth, a disagreement being asked about once and otherwise flagged.

\paragraph{Ingestion.}
The current ingestor stores chunks verbatim with their vectors (a table is one chunk, whole, never cut), tags prose sentences with their section path and the concept captions they mention, writes the document profile, and extracts nothing. Ingesting a 10-K costs the two calls above plus embeddings, about 30 seconds and \$0.15--0.30. (The earlier ingestor that read every table into facts through a model is kept as version~1; it costs about \$2--3 per filing and is the one whose decisions Table~\ref{tab:commitment} measures.)

\subsection{The source graph and the bridge}
\label{sec:impl-source}

The source graph (Section~\ref{sec:idea-three}) is authored as JSON and built into a database by a builder that refuses the whole build on a quote not printed in its official text, an unknown period word, a dangling relation or an unbacked ``measured'' entry. Nineteen official texts were fetched under a manifest (the forms' instructions, Regulations S-K and S-X, the EDGAR form list, the dei dictionary). The closed period vocabulary was checked before adoption: it names 98.0\% of the values in 10-Ks, 99.3\% in 10-Qs and 73\% in S-1 registration statements; the quarter rule agrees with the SEC's own quarter code on 99.2\% of 21,692 10-Qs. The graph is consulted at three points: the cover reading (candidate kinds), the document search (which stored document covers the asked period: a prior year held as a comparative ranks below the document's own year and says so; a Q4 question fits no 10-Q and gets the graph's note), and the outline step (the skeleton hint and part linking).

The bridge is measured by one script over the SEC Financial Statement and Notes data sets: for each of 32,288 filings of seven kinds, every consolidated value is placed by its statement code and its period word relative to the filing's own period, and grouped into rows of (concept, kind, part, location, periods, count, of). Using the Notes data set rather than the plain quarterly files was decisive: the plain files hold only the face of the statements, which is why an earlier count found the effective tax rate in 14 filings; the bridge finds it in the income-tax note of 56\% of 10-Ks. The bridge gives the question-answering loop the \emph{chance} a document of this kind, from a filer in this industry, prints the target (searched before any formula when high, derived first when measured low, searched first when unknown) and the \emph{parts} to search first. Every accepted find is written back as an observed row.

\subsection{Retrieval}
\label{sec:impl-retrieval}

\paragraph{The ladder (version 2).} The control retriever runs fixed steps in a fixed order: for a number, a table index first (every stored table parsed just enough to find and cite it: header rows, row labels, every numeric cell with its row and column; scored by row-label coverage, idf-weighted word overlap, heading match and whether the coordinate's words are printed near the table), then prose routed by heading, by words, by meaning; tables come back whole with the chunks before and after as context.

\paragraph{The learned scorer (version 3).}
The ladder consulted four kinds of evidence one at a time with hand-picked weights. The replacement is \emph{one score per chunk} from thirteen signals (words, synonym, table row label, heading path, meaning by vector, the bridge's share for the part in force, the part's content tag, the coordinate's words near and far, table, table$\times$numeric, numbers printed$\times$numeric, length), fitted as McFadden's conditional logit~\citep{mcfadden1974} with a ``not printed in this document'' alternative, so that it returns a probability over the chunks \emph{and} a probability that the target is not there. The data arrangement is the one principle of Section~\ref{sec:idea} applied: 420 filings sampled from two SEC quarters disjoint from the bridge's five, labelled by their own inline-XBRL tags (178,185 values and 23,195 text blocks with the chunks that print them, no model call and no benchmark question), the FinanceBench companies excluded by CIK, held out by company inside the fit. The served probability mass and the bridge's counts are then calibrated against held-out filings by isotonic regression~\citep{zadrozny2002}, because both were found to under-predict in the middle of their range (a raw served mass of 0.09 was a hit 26\% of the time). Version~3 is the registry's default (Figure~\ref{fig:retrieval}).

\begin{figure}[t]
\centering
\fitwidth{\begin{tikzpicture}[
  >=Stealth, font=\scriptsize,
  box/.style={draw,thick,rounded corners=3pt,fill=cLight,align=center,inner sep=4pt,text width=30mm},
  sig/.style={draw=cCode,thin,rounded corners=2pt,fill=white,align=left,inner sep=3pt,text width=38mm,font=\tiny},
  arr/.style={->,thick,gray!80}
]
\node[box,draw=cDoc] (need) at (0,0) {\textbf{need}\\ words, synonyms (filer captions of the mapped entity), coordinate, numeric or text, scope (parts), concept};
\node[sig,right=8mm of need,yshift=12mm] (s1) {\textbf{wording}: words in chunk $\cdot$ synonym $\cdot$ table row label $\cdot$ heading path};
\node[sig,below=2mm of s1] (s2) {\textbf{meaning}: cosine of need and chunk vectors};
\node[sig,below=2mm of s2] (s3) {\textbf{place} (from the bridge and source graph): share printed in the part in force $\cdot$ part content tag};
\node[sig,below=2mm of s3] (s4) {\textbf{location words}: company and year printed near / far};
\node[sig,below=2mm of s4] (s5) {\textbf{form}: table $\cdot$ table$\times$numeric $\cdot$ numbers printed$\times$numeric $\cdot$ length};
\node[box,draw=cCode,right=8mm of s3,text width=34mm] (logit) {\textbf{conditional logit}\\ $P(\text{chunk } c)=\dfrac{e^{w\cdot x_c}}{\sum_{c'} e^{w\cdot x_{c'}}+e^{u}}$\\[2pt] one extra alternative $u$: ``not printed in this document''};
\node[box,draw=cCode,right=8mm of logit,text width=30mm] (cal) {\textbf{calibration}\\ isotonic map of served mass $\to$ observed hit rate, fitted on held-out companies};
\node[box,draw=cDoc,below=6mm of cal,text width=30mm] (out) {\textbf{passages}, most probable first; tables whole with context; per need: mass served, unserved, not in document};
\draw[arr] (need.east) -- ++(4mm,0) |- (s1.west);
\draw[arr] (need.east) -- ++(4mm,0) |- (s5.west);
\foreach \s in {s1,s2,s3,s4,s5} {\draw[arr] (\s.east) -- ++(4mm,0) |- (logit.west);}
\draw[arr] (logit) -- (cal); \draw[arr] (cal) -- (out);
\node[gray,align=left,font=\tiny,below=4mm of s5.south west,anchor=north west,text width=100mm] {Fitted on 420 EDGAR filings from two SEC quarters disjoint from the bridge's, labelled by their own inline-XBRL tags (178k values, 23k text blocks), FinanceBench companies excluded by CIK, held out by company. Largest standardised weights: numbers printed $+1.15$, meaning $+1.00$, synonym $+0.41$, words $+0.40$, row $+0.30$.};
\end{tikzpicture}}
\caption{The learned retriever (version 3). Thirteen signals per chunk, grouped by the kind of evidence they carry, enter one conditional logit with a ``not printed'' alternative; the served probability is then calibrated. The bridge and the source graph enter as signals rather than as a fixed step.}
\label{fig:retrieval}
\end{figure}


\subsection{The no-graph control}
\label{sec:impl-nograph}

Before the graph was put back, a complete no-graph stack was built as the control and the base to build on: an ingestor that stores text and tables only, the ladder retriever above, and a loop that plans needs, retrieves, reads tables whole with their context and grounds every figure. Its grounding rules are the ones the current loop still uses: a table figure is cited as passage, quote, row and unit words, and the code checks the cell is printed in that row of that table, takes the value and sign from the cell, accepts the unit when its quote is found in the heading or context, records a unit claimed without a quote as an assumption, and evaluates arithmetic itself with base-unit normalisation and dimensional checking (cross-family addition refused, like-family ratios dimensionless). A prose figure's value is read from the model's quote, which must be printed in the passage, with brackets and minus signs read as negative.

\subsection{The question-answering loop (version 4)}
\label{sec:impl-loop}

The current loop puts the graph back as the prior (Figure~\ref{fig:loop}). Its pieces are plain data hand-offs; a model is called at eight points only, each with one narrow choice, and everything between them is code.

\begin{figure}[p]
\centering
\fitheight{\fitwidth{\begin{tikzpicture}[
  >=Stealth, font=\scriptsize, node distance=3.5mm,
  step/.style={draw,thick,rounded corners=3pt,align=center,inner sep=4pt,text width=52mm,minimum height=11mm},
  model/.style={step,draw=cModel,fill=cModel!8},
  code/.style={step,draw=cCode,fill=cLight},
  kbnote/.style={draw=cConnector,thick,rounded corners=2pt,fill=cConnector!10,align=left,inner sep=3pt,text width=48mm,font=\tiny},
  arr/.style={->,thick,gray!80},
  back/.style={->,thick,cModel,dashed}
]
\node[model] (und) at (0,0) {\textbf{1 Understand} \textcolor{cModel}{[model]}\\ targets + coordinates, formula over targets, givens, unit, analysis statement};
\node[model,below=of und] (map) {\textbf{2 Map words $\to$ entity} \textcolor{cModel}{[model]}\\ candidates by words, filer captions and meaning; a number or \emph{none}; never cached};
\node[model,below=of map] (docs) {\textbf{3 Choose documents} \textcolor{cModel}{[model, or code when one fits]}\\ tools: search documents, company filers, list};
\node[code,below=of docs] (dec) {\textbf{4 Search or derive?} [code]\\ calibrated chance from the bridge, nearest industry first:\\ high $\to$ search; measured 0 $\to$ formulas; unknown $\to$ search then derive};
\node[model,below=of dec] (routes) {\textbf{5 Routes} \textcolor{cModel}{[model when several]}\\ formulas solved for the target, ranked by trust then joint chance of inputs; ancestors excluded; a definition stated by the model when the graph has none};
\node[model,below=of routes] (find) {\textbf{6 Find} \textcolor{cModel}{[model, rounds $\le 3$]}\\ scope: the parts the bridge names; passages from the retriever; reply: found / found text / collect / widen / search / absent};
\node[code,below=of find] (ground) {\textbf{Ground} [code]\\ quote printed in the passage and row; value and sign from the cell; unit from quoted evidence; kind of number vs kind of quantity};
\node[code,below=of ground] (plan) {\textbf{7 Plan state} [code]\\ settle, propagate up complete routes; switch to a cheaper paused route; abandon inputs no longer needed; evidence kept in full};
\node[model,below=of plan] (resp) {\textbf{8 Respond} \textcolor{cModel}{[model; none for a bare number]}\\ every derived value recomputed from cited inputs via the graph's equation; claims must quote evidence; may ask for more targets};
\node[code,below=of resp] (eval) {\textbf{Evaluate} [code; model only to restate]\\ shape, unit, sign fixed; percentage points for rate differences};
\draw[arr] (und) -- (map); \draw[arr] (map) -- (docs); \draw[arr] (docs) -- (dec); \draw[arr] (dec) -- (routes);
\draw[arr] (dec.east) -- ++(6mm,0) |- node[pos=.25,right,font=\tiny,align=left]{search\\first} (find.east);
\draw[arr] (routes) -- (find); \draw[arr] (find) -- (ground); \draw[arr] (ground) -- (plan); \draw[arr] (plan) -- (resp); \draw[arr] (resp) -- (eval);
\draw[back] (ground.west) -- ++(-5mm,0) |- node[pos=.25,left,font=\tiny,align=right]{rejected:\\ re-ask once} (find.west);
\draw[back] (plan.west) -- ++(-10mm,0) |- node[pos=.3,left,font=\tiny,align=right]{input absent:\\ expand through\\ a formula, or\\ switch route} (routes.west);
\draw[back] (resp.west) -- ++(-15mm,0) |- node[pos=.3,left,font=\tiny,align=right]{needs:\\ new targets\\ ($\le 2$ rounds)} (map.west);

% knowledge-base notes on the right
\node[kbnote,right=12mm of map] (n1) {\textcolor{cEntity}{\textbf{concept graph}}: names, aliases, SEC concept links, filer captions, vectors};
\node[kbnote,right=12mm of docs] (n2) {\textcolor{cDoc}{\textbf{source graph}} + company store: which stored document covers the company and period, as its own year or a comparative};
\node[kbnote,right=12mm of dec] (n3) {\textcolor{cConnector}{\textbf{bridge}}: share of filings of this kind and industry printing the entity; in which part and statement};
\node[kbnote,right=12mm of routes] (n4) {\textcolor{cEntity}{\textbf{concept graph}}: equations and subtotals that produce the entity, solved symbolically; trust: identity $<$ law $<$ empirical};
\node[kbnote,right=12mm of find] (n5) {\textcolor{cConnector}{\textbf{bridge}} $\to$ parts to search first; \textcolor{cDoc}{\textbf{source graph}} $\to$ headings linked to parts; retriever v3 ranks chunks};
\node[kbnote,right=12mm of resp] (n6) {\textcolor{cEntity}{\textbf{concept graph}}: the equation reloaded by connector id for recomputation in code};
\foreach \a/\b in {n1/map,n2/docs,n3/dec,n4/routes,n5/find,n6/resp} {\draw[arr,gray!50] (\a.west) -- (\b.east);}
\end{tikzpicture}}}
\caption{The question-answering loop (version 4). Plum boxes are the points where a model is called, each with one narrow choice; grey boxes are code. Dashed arrows are the loop's backtracking paths. The notes on the right say which knowledge base each step consults.}
\label{fig:loop}
\end{figure}


\begin{enumerate}
\item \textbf{Understand.} One call over the question alone returns the targets as the question names them, each with its coordinate (company, period, segment), a formula over target ids only when the question spells out a change, growth, average or sum, the given values, the unit asked, and an \emph{analysis statement} when something is to be judged (``is liquidity looking good'' becomes current and quick ratios at two year ends plus a text target for what management says, with the judgement left to the end). A named quantity such as ``quick ratio'' is one target; the graph decides its parts.
\item \textbf{Map.} Each target's words go to one entity or none: candidates by the words (an entity whose name or alias is the words; every entity whose SEC concept filers label with the words, with a floor on filings) and by meaning (nearest vectors), then one call showing the question, the words and each candidate's definition and printed captions. Nothing is cached. A rejected target is searched by text alone.
\item \textbf{Choose documents.} Per coordinate group, the model calls tools (search the stored documents by company, period and kind; the company store of filers and co-registrants; list all) and selects only sources that appeared. When exactly one document fits every key on its own cover the code chooses with no call. The company store (16,416 filers, 62,943 names over time, 548 co-registrant pairs built from the SEC's filing indexes) answers ``whose documents may hold this company's numbers'', so Entergy Texas finds the Entergy combined 10-K.
\item \textbf{Decide: search or derive.} The bridge's calibrated chance that a document of this kind prints the entity, counted nearest-first (company, four-digit industry, two-digit, division, all). High: search. Measured zero: formulas first. Unknown (no concept link, or a kind with no filings): search, and derive on a miss.
\item \textbf{Routes.} The relations that produce an entity, solved symbolically into readable formulas, ranked by trust then by the joint chance of their inputs (counted in sequence, least likely first, because findings are strongly linked: current liabilities is reported by 98.8\% of filers reporting current assets and 3.9\% of the others); one usable route is taken by code, several go to one call that may decline. A target the graph does not know and the document does not print may have its usual definition stated by the model once, over simpler quantities, kept as an assumption. Formulas using an ancestor of the target are excluded, which stops ``short-term investments'' being derived from the quick ratio it feeds.
\item \textbf{Find.} For each chosen document, in the parts the bridge names first and then the whole document: the retriever's most probable chunks, shown as whole tables with context or whole prose chunks, together with the question, the assumptions so far, the values already found and what has been learned about this document; the model replies \emph{found} (passage, quote, row, unit, unit quote), \emph{found text}, \emph{collect} (passages for a text target), \emph{widen}, \emph{search} with other words, or \emph{absent}. Every citation is checked as in Section~\ref{sec:impl-nograph}; a mismatch between the kind of number cited and the kind of quantity sought (currency for a ratio) is sent back once. The passage and a one-line ``learned'' note are returned as evidence.
\item \textbf{Plan and backtrack.} The planner keeps the whole state as memory: targets with status, groups with the tool transcript, routes with inputs and expected remaining cost, evidence in full. A found input settles and propagates up every route whose inputs are complete; a failed route activates the cheapest paused one; a derived target closes its other routes and abandons their inputs. Mapping, document choice and each round of searches run in parallel within a question.
\item \textbf{Respond and evaluate.} Every derived value is recomputed in code from its cited inputs through the equation reloaded from the graph; a numeric question asking for a number and nothing else is answered with no model call. Otherwise one call sees every figure with how it was obtained, every passage verbatim, the notes and assumptions, and answers in a fixed shape whose claims must quote the evidence and whose numbers must match a figure; a failed reply is sent back once. The evaluator then fixes shape, unit and sign in code (a ratio asked in percent, a decrease asked as a magnitude, percentage points for a difference of rates), calling a model only to restate a verdict the code cannot read off.
\end{enumerate}

\subsection{Tooling}

Three tools were built to see what the system does. A \emph{trace viewer} renders a run as a conversation between the pieces of the loop, one column per piece, every model call a box with a one-line summary of its reply and the full exchange and planner state a click away. A \emph{knowledge explorer} shows the three knowledge bases as one page with one item in focus and its neighbours in named columns (a quantity's formulas on one side and where it is printed on the other). And two drivers run the benchmarks end to end with every call logged, resumable, several questions at a time. All of them were used in the analyses of Section~\ref{sec:results}.
